% =============================================================================
% Paper 1 — Present-epoch scale matching and C_obs hygiene on cubic T^3
% Genre: technical note (no-gos + one derived EFT identity)
% Class: Physical Review D (revtex4-2)
% =============================================================================
\documentclass[aps,prd,reprint,superscriptaddress,nofootinbib,floatfix]{revtex4-2}

\usepackage[T1]{fontenc}
\usepackage[utf8]{inputenc}
\usepackage{amsmath,amssymb,amsthm}
\usepackage{booktabs}
\usepackage{graphicx}
\usepackage[hidelinks]{hyperref}

\graphicspath{{figures/}{../../Assets/Figures/}{../../Assets/Archive/}}

\newcommand{\Tthree}{T^{3}}
\newcommand{\Ttwo}{T^{2}}
\newcommand{\Cobs}{C_{\mathrm{obs}}}
\newcommand{\Cm}{C_{m}}
\newcommand{\CIR}{C_{\mathrm{IR}}}
\newcommand{\azero}{a_{0}}
\newcommand{\Hzero}{H_{0}}

\begin{document}

\title{Present-epoch scale matching, the weak-field invariant
\texorpdfstring{\(\Cobs\)}{C\_obs}, and no-go results for a cubic
\texorpdfstring{\(\Tthree\)}{T3} geometric story}

\author{Brendon Boyd}
\email{brendon.boyd@itsm-cosmology.org}
\affiliation{Independent Researcher, Burswood, Western Australia, Australia}
\thanks{%
\begin{tabular}{@{}l@{\ }l@{}}
ORCID: & \href{https://orcid.org/0009-0007-4177-2612}{0009-0007-4177-2612} \\
GitHub: & \href{https://github.com/brendohxd/ITSM-Integrated-Toroidal-Syntropic-Model}{\texttt{github.com/brendohxd/ITSM-...}} \\
Zenodo: & \href{https://doi.org/10.5281/zenodo.18808348}{doi:10.5281/zenodo.18808348} \\
\end{tabular}%
}

\date{\today}

\begin{abstract}
The coincidence \(2\pi\azero\sim c\Hzero\) is a standard observation in the
MOND literature.  Several recent proposals attempt to elevate it to a
\emph{geometric derivation} by combining (i)~a compact flat cubic three-torus
of present-day side \(L=c/\Hzero\) with a putative ``circulation quantum''
\(\kappa=cL\), and (ii)~the isotropic projector ratio
\(\operatorname{Tr}(h)/\operatorname{Tr}(\gamma)=2/3\) with the weak-field
coefficient in an AQUAL-type action.  We show that neither step is presently
a derivation.  Writing \(\kappa=cL\) does not quantize circulation; the package
\(L=c/\Hzero\), \(\Gamma=cL\), and \(\omega=\Hzero\) is kinematically
inconsistent if \(L\) is a circumference; and cubic \(E_1\) topology bounds
require \(L\gtrsim\mathcal{O}(6)\,c/\Hzero\), excluding the \(\alpha=1\)
matching that reproduces the MOND scale.  Exact tracking
\(L_{\mathrm{phys}}(t)=c/H(t)\) is incompatible with fixed comoving moduli in
ordinary FLRW except for a coasting expansion.  The projector identity is
correct but topology-independent and does not appear in the weak-field action,
so it cannot fix
\(\Cobs=\Cm^{3/2}/\sqrt{\CIR}\), which we derive as the unique invariant of
field rescaling.  Imposing both \(\azero=c\Hzero/2\pi\) and \(\Cobs=2/3\)
yields \(a_{0,\mathrm{eff}}=(4/9)\azero\), in conflict with the empirical RAR
normalization.  What remains is a present-epoch phenomenological encoding of
the coincidence, a clean EFT invariant, and an explicit reconstruction checklist
for any future geometric claim.
\end{abstract}

\keywords{Modified Newtonian dynamics;
cosmic topology;
effective field theory;
radial acceleration relation;
compact manifolds}

\maketitle

% =============================================================================
\section{Introduction}
\label{sec:intro}

Empirical galactic dynamics exhibit a characteristic acceleration
\(\azero\sim 1.2\times 10^{-10}\,\mathrm{m\,s^{-2}}\) that organizes the
radial acceleration relation (RAR) and the baryonic Tully--Fisher relation
\cite{mond,famaey_mcgaugh,mcgaugh_rar,lelli_rar}.
The numerical coincidence
\begin{equation}
  2\pi\,\azero \sim c\,\Hzero
  \label{eq:coincidence}
\end{equation}
is well known \cite{mond,famaey_mcgaugh}.  It is a legitimate research question
whether a compact flat three-torus \(\Tthree=S^{1}\times S^{1}\times S^{1}\)
\cite{lachieze_rey_luminet,planck_topology} and a weak-field effective field
theory (EFT) can turn that coincidence into a derivation.

A recurring two-step proposal answers ``yes'':
\begin{enumerate}
\item[(A)] ``Circulation quantization'' on a domain of size \(L=c/\Hzero\)
  yields \(\azero=c\Hzero/2\pi\);
\item[(B)] the transverse-projector ratio \(2/3\) yields the EFT coefficient
  \(\Cobs\) in a deep-MOND square-root law.
\end{enumerate}
Both steps fail under standard technical scrutiny.  This paper is a short
technical note: it withdraws those derivations, records the associated no-go
results, and isolates what \emph{does} follow cleanly from a static AQUAL-type
action \cite{bekenstein_milgrom}.  Table~\ref{tab:status} summarizes the
status of each claim.

% Full-width, non-floating table: spans both columns; following text starts after it.
\onecolumngrid
\vspace{0.9em}
\refstepcounter{table}\label{tab:status}
\noindent
\begin{minipage}{\textwidth}
\small
\noindent\textbf{TABLE~\thetable.\ }
Status of claims discussed in this work.\\[0.55em]
\begin{ruledtabular}
\begin{tabular}{@{}p{0.50\textwidth}@{\hspace{0.8em}}p{0.12\textwidth}@{\hspace{0.8em}}p{0.30\textwidth}@{}}
\textbf{Quantity / claim} & \textbf{Status} & \textbf{Note} \\
\colrule
Empirical coincidence \(2\pi\azero\sim c\Hzero\)
  & Fact & Motivating observation only \\
\hline
\(\azero\equiv c\Hzero/2\pi\) at the present epoch
  & Phen. & Named postulate; not derived \\
\hline
``Circulation quantization \(\Rightarrow\azero\)''
  & Withdrawn & \(\kappa=cL\) is not \(h/m\) quantization \\
\hline
Package \((L,\Gamma=cL,\omega=\Hzero)\) as prior proposals
  & Withdrawn & Inconsistent if \(L\) is a circumference \\
\hline
Cubic \(E_1\) with \(L=c/\Hzero\) vs.\ Planck topology
  & Excluded & \(L\gtrsim\mathcal{O}(6)\,c/\Hzero\) required \\
\hline
\(L_{\mathrm{phys}}(t)=c/H(t)\) with fixed comoving moduli
  & Excluded & Needs evolving moduli or coasting FLRW \\
\hline
Isotropic count \(\operatorname{Tr}(h)/\operatorname{Tr}(\gamma)=2/3\)
  & Derived & Generic 3D identity; not special to \(\Tthree\) \\
\hline
\(\Cobs=\Cm^{3/2}/\sqrt{\CIR}\)
  & Derived & Field-normalization invariant \\
\hline
\(\Cobs=2/3\) from the trace ratio
  & Withdrawn & Projector absent from the action \\
\hline
Simultaneous \(\azero=c\Hzero/2\pi\) and \(\Cobs=2/3\) as RAR fit
  & Excluded & Predicts \(a_{0,\mathrm{eff}}=(4/9)\azero\) \\
\hline
Compensated sources on compact \(\Tthree\)
  & Derived & \(\int\mathrm{div}=0\) \\
\hline
Geometric derivation of \(\azero\) or \(\Cobs\)
  & Open & Path in Sec.~\ref{sec:path} \\
\end{tabular}
\end{ruledtabular}
\end{minipage}
\vspace{1.0em}
\twocolumngrid

% =============================================================================
\section{Cubic three-torus geometry and notation}
\label{sec:t3}

A flat three-torus is the quotient \(\mathbb{R}^{3}/\Lambda\) by a lattice of
translations \cite{lachieze_rey_luminet}.  The correct schematic is a
parallelepiped with opposite faces identified (Fig.~\ref{fig:t3}), not the
doughnut surface of revolution in \(\mathbb{R}^{3}\).  That doughnut is an
embedding of \(\Ttwo\).

\begin{figure}[t]
\centering
\includegraphics[width=\columnwidth]{itsm_t3_fundamental_domain.pdf}
\caption{Flat cubic \(\Tthree\) fundamental domain: opposite faces identified
(\(A\sim A'\), \(B\sim B'\), \(C\sim C'\)).  Edge lengths \(L_x,L_y,L_z\) are
the three independent periods.  Cubic means \(L_x=L_y=L_z\).  This is not the
embedded doughnut \(\Ttwo\).}
\label{fig:t3}
\end{figure}

``\(\Tthree\)'' alone does not imply equal periods.  Isotropic statements
require a \emph{cubic} (or explicitly isotropised) lattice.  Rectangular and
twisted flat manifolds (\(E_2\), \(E_3\), \ldots) are different geometries with
different CMB constraints \cite{planck_topology,compact_promise}.

Throughout, \(L\) means a \emph{physical} edge length at the epoch under
discussion unless ``comoving'' is written explicitly.  When a numerical Hubble
scale is needed we take
\(\Hzero=70\,\mathrm{km\,s^{-1}\,Mpc^{-1}}\) for illustration; none of the
no-go arguments depends on that choice at the factor-of-two level.

% =============================================================================
\section{\texorpdfstring{\(\azero=c\Hzero/2\pi\)}{a0 = c H0 / 2 pi}
is not circulation quantization}
\label{sec:a0}

\subsection{What quantization requires}

In a superfluid with order parameter \(\Psi=\sqrt{n}\,e^{i\theta}\), the
velocity follows from the phase gradient and
\begin{equation}
  \Gamma = \oint \mathbf{v}\cdot d\mathbf{l}
  = n\,\frac{h}{m}\,,
  \qquad n\in\mathbb{Z}\,,
  \label{eq:true-quant}
\end{equation}
with single-valuedness of \(\Psi\) on closed cycles
\cite{onsager,feynman,volovik}.  Equation~\eqref{eq:true-quant} quantizes
\emph{phase winding}, not ``length times a speed limit.''

\subsection{What the geometric proposals actually write}

Those proposals introduce
\begin{equation}
  \kappa_{\mathrm{cosmo}} \stackrel{?}{=} c\,L
  \label{eq:fake-kappa}
\end{equation}
by multiplying a maximum speed by a length.  For constant speed on a cycle of
length \(L\), one has only \(\Gamma=vL\); setting \(v=c\) yields \(\Gamma=cL\).
No integer winding, no order parameter, and no mass \(m\) enter.  That is a
kinematic assignment, not quantization.

If one \emph{did} set the \(n=1\) quantum equal to \(c^{2}/\Hzero\),
\begin{equation}
  m = \frac{h\Hzero}{c^{2}}
  \sim 10^{-32}\,\mathrm{eV}/c^{2}\,,
  \label{eq:mass-pred}
\end{equation}
the microscopic mass reappears as a sharp prediction---the opposite of claiming
that no microphysics is needed.

\subsection{Internal inconsistency of the \texorpdfstring{\((L,\Gamma,\omega)\)}{(L, Gamma, omega)} package}

Suppose \(L=c/\Hzero\) is the physical circumference of a closed cycle, so
\(R=L/(2\pi)\).  Then:
\begin{itemize}
\item \(\Gamma=cL\) implies \(v=c\);
\item \(\omega=\Hzero\) with \(v=\omega R\) implies
  \(v=\Hzero L/(2\pi)=c/(2\pi)\), contradicting \(v=c\);
\item the centripetal acceleration \(c^{2}/R\) is \(2\pi c\Hzero\), not
  \(c\Hzero/2\pi\).
\end{itemize}
Thus one cannot simultaneously maintain
\begin{equation}
  L=\frac{c}{\Hzero}\,,
  \qquad
  \Gamma=cL\,,
  \qquad
  \omega=\Hzero
  \label{eq:inconsistent}
\end{equation}
under that geometric reading of \(L\).  ``Distributing \(c\Hzero\) across a
\(2\pi\) angular topology'' is not coordinate-invariant: the range of an
angular chart can be rescaled.  A physical \(2\pi\) arises only from a defined
phase convention coupled to velocity.

\paragraph{Conclusion.}
\begin{equation}
  \azero \equiv \frac{c\Hzero}{2\pi}
  \label{eq:a0-phen}
\end{equation}
may be retained only as a \emph{present-epoch phenomenological scale-matching
postulate} that encodes Eq.~\eqref{eq:coincidence}.  For
\(\Hzero=70\,\mathrm{km\,s^{-1}\,Mpc^{-1}}\) one finds
\(\azero\approx 1.1\times 10^{-10}\,\mathrm{m\,s^{-2}}\), of the observed
order of magnitude \cite{mcgaugh_rar}.  All ``geometric derivation from
circulation quantization'' language should be withdrawn.

% =============================================================================
\section{Cubic topology bounds exclude \texorpdfstring{\(L=c/\Hzero\)}{L = c/H0}}
\label{sec:bounds}

Planck topology analyses constrain cubic \(E_1\) three-tori tightly.  In terms
of the radius \(R_i\) of the largest sphere inscribed in the fundamental
domain, Planck reports \(R_i>0.97\,\chi_{\mathrm{rec}}\) for the cubic case at
high confidence \cite{planck_topology}.  For a cube of side \(L\),
\(R_i=L/2\), so
\begin{equation}
  L \gtrsim 1.94\,\chi_{\mathrm{rec}}\,.
  \label{eq:L-chi}
\end{equation}
With \(\chi_{\mathrm{rec}}\simeq 3.15\,c/\Hzero\) in a Planck cosmology
\cite{planck},
\begin{equation}
  L \gtrsim 6\,\frac{c}{\Hzero}
  \label{eq:L-H}
\end{equation}
as an order-of-magnitude statement (exact numerical factors depend on the
statistic and priors).  Recent topology surveys likewise place the shortest
\(E_1\) side near the diameter scale of the last-scattering surface, not one
Hubble radius \cite{compact_promise,planck_topology}.

Write \(L=\alpha\,c/\Hzero\).  Empirical scale matching uses \(\alpha=1\);
cubic bounds require \(\alpha\sim\mathcal{O}(6)\).  If the phenomenological
formula is reinterpreted as \(\azero=c^{2}/(2\pi L)=c\Hzero/(2\pi\alpha)\),
topology-safe \(\alpha\) yields
\begin{equation}
  \azero \lesssim \frac{1}{6}\times\frac{c\Hzero}{2\pi}
  \sim 1.8\times 10^{-11}\,\mathrm{m\,s^{-2}}\,,
  \label{eq:a0-suppressed}
\end{equation}
destroying the intended MOND-scale match.  Therefore the cubic-\(E_1\),
\(\alpha=1\) package that simultaneously targets \(\azero\sim 10^{-10}\) and
CMB non-detection is a no-go.

Vague statements that a ``Hubble-scale domain is compatible with non-detection''
are incorrect for cubic \(E_1\).

% =============================================================================
\section{Present-epoch matching versus fixed-comoving FLRW}
\label{sec:dynamic}

For fixed comoving moduli \(L^{\mathrm{(com)}}\),
\begin{equation}
  L_{\mathrm{phys}}(t)=a(t)\,L^{\mathrm{(com)}}\,,
  \qquad
  \frac{\dot L_{\mathrm{phys}}}{L_{\mathrm{phys}}}=H\,.
  \label{eq:fixed}
\end{equation}
Requiring instead \(L_{\mathrm{phys}}(t)=c/H(t)\) gives
\begin{equation}
  \frac{\dot L_{\mathrm{phys}}}{L_{\mathrm{phys}}}
  = -\frac{\dot H}{H}\,,
  \label{eq:dsm-rate}
\end{equation}
and equating the two yields \(\dot H=-H^{2}\), i.e.\ deceleration parameter
\(q=0\) (coasting).  Exact tracking is therefore incompatible with ordinary
matter-, radiation-, or \(\Lambda\)-dominated FLRW evolution under fixed
comoving moduli.

The remaining alternatives are: (i)~present-epoch numerical identification
only; or (ii)~an independent modulus sector with its own kinetic term and
stress-energy.  No such sector is constructed here.

The Hubble sphere \(c/H\) is also not generally a particle or event horizon;
calling it a ``global causal domain'' is physically misleading
\cite{davis_lineweaver_horizon}.

% =============================================================================
\section{The projector ratio does not determine
\texorpdfstring{\(\Cobs\)}{C\_obs}}
\label{sec:proj}

\subsection{The identity}

For any unit \(n^{i}\) in three-dimensional Riemannian geometry,
\begin{equation}
  h_{ij}=\gamma_{ij}-n_{i}n_{j}\,,
  \qquad
  \frac{\operatorname{Tr}(h)}{\operatorname{Tr}(\gamma)}=\frac{2}{3}\,.
  \label{eq:trace}
\end{equation}
This is correct, topology-independent, and independent of superfluidity.  It
discards tensorial orientation information and is not special to \(\Tthree\).

\subsection{Static AQUAL-type action}

A static AQUAL-type force sector \cite{bekenstein_milgrom} may be written
\begin{align}
  S_{\mathrm{WF}}
  &= \int\! dt\,d^{3}x
  \Biggl[
  -\frac{|\nabla\Phi_{N}|^{2}}{8\pi G}
  - \rho_{b}\Phi_{N}
  \nonumber\\
  &\qquad
  - \frac{\CIR}{12\pi G\,\azero}\,|\nabla\psi|^{3}
  - \Cm\,\rho_{b}\,\psi
  \Biggr]\,,
  \label{eq:SWF}
\end{align}
with \(q=|\nabla\psi|=\sqrt{\gamma^{ij}\partial_{i}\psi\partial_{j}\psi}\).
Neither \(n^{i}\) nor \(h_{ij}\) appears.  The kinetic operator is fully
three-dimensional and isotropic.  There is therefore no field-theoretic sense
in which Eq.~\eqref{eq:trace} ``matches'' onto the EFT.

Variation yields
\begin{equation}
  \nabla\cdot\!\left(\frac{|\nabla\psi|}{\azero}\,\nabla\psi\right)
  = 4\pi G\,\frac{\Cm}{\CIR}\,\rho_{b}\,.
  \label{eq:p3}
\end{equation}
In spherical symmetry, with \(g_{N}=|\nabla\Phi_{N}|\) and physical
acceleration \(g_{P}=\Cm|\nabla\psi|\),
\begin{equation}
  g_{P} = \Cobs\sqrt{\azero\,g_{N}}\,,
  \qquad
  \boxed{
  \Cobs = \frac{\Cm^{3/2}}{\sqrt{\CIR}}
  }\,.
  \label{eq:Cobs}
\end{equation}

\subsection{Field-rescaling invariance}

Under \(\psi\mapsto s\psi\) with \(s>0\),
\begin{equation}
  \Cm\mapsto\frac{\Cm}{s}\,,
  \qquad
  \CIR\mapsto\frac{\CIR}{s^{3}}\,,
  \label{eq:rescale}
\end{equation}
while \(\Cobs\) is invariant.  Thus \(\Cobs\) is the unique combination of the
two positive coefficients that can enter observables.  Setting \(\Cm=\CIR\) so
that \(\Cobs=C\) is a normalization convention, not a calculation of
\(C=2/3\).

\begin{figure}[t]
\centering
\includegraphics[width=\columnwidth]{itsm_23_factor_schematic.pdf}
\caption{Isotropic degree-of-freedom cartoon only.  It does not appear in
Eq.~\eqref{eq:SWF} and must not be read as a derivation of \(\Cobs\).}
\label{fig:cproj}
\end{figure}

\paragraph{Withdrawn.}
Any statement that \(\Cobs=2/3\) follows from \(\Tthree\) geometry, from
Eq.~\eqref{eq:trace}, or from ``matching'' without inserting \(h^{ij}\) into
the action and integrating out fields.

\subsection{Compact sources}

On compact \(\Tthree\), the integral of any pure divergence vanishes, so
localized solutions must use a compensated density contrast (and, if present,
a zero-mean exchange source).  For general morphology,
\begin{equation}
  \frac{|\nabla\varphi|}{\azero}\,\nabla\varphi
  = \frac{\Cm}{\CIR}\,\nabla\Phi_{N} + \nabla\times\mathbf{H}\,,
  \label{eq:curl}
\end{equation}
so pointwise square-root formulae are approximations for disks.  These
statements remain valid; they do not resurrect \(\Cobs=2/3\).

% =============================================================================
\section{Mutual inconsistency with the empirical RAR scale}
\label{sec:rar}

In the deep-MOND regime, \(g_{P}=\Cobs\sqrt{\azero g_{N}}\) is equivalent to
\begin{equation}
  g_{P}=\sqrt{a_{0,\mathrm{eff}}\,g_{N}}\,,
  \qquad
  a_{0,\mathrm{eff}} = \Cobs^{2}\,\azero\,.
  \label{eq:a0eff}
\end{equation}
Imposing both \(\azero=c\Hzero/2\pi\approx 1.1\times 10^{-10}\,\mathrm{m\,s^{-2}}\)
and \(\Cobs=2/3\) yields
\begin{equation}
  a_{0,\mathrm{eff}} = \frac{4}{9}\azero
  \approx 4.9\times 10^{-11}\,\mathrm{m\,s^{-2}}\,,
  \label{eq:a0eff-num}
\end{equation}
whereas modern RAR determinations find
\(a_0\simeq 1.2\times 10^{-10}\,\mathrm{m\,s^{-2}}\)
\cite{mcgaugh_rar,lelli_rar}, a factor \(\sim 2.4\) discrepancy.
Equivalently, asymptotic baryonic Tully--Fisher velocities scale as
\(v^{4}\propto \Cobs^{2}\azero GM_{b}\); \(\Cobs=2/3\) lowers \(v\) by
\(\sim 18\%\) at fixed baryonic mass, which is observationally large.

Therefore one cannot simultaneously advertise
\(\azero=c\Hzero/2\pi\) as the empirical MOND scale and impose a universal
\(\Cobs=2/3\).  Consistent options for any future fit are:
\begin{enumerate}
\item set \(\Cobs=1\) (up to the definition of \(\azero\));
\item redefine the input scale so that \(\Cobs^{2}\azero\) matches the RAR;
\item drop universality and derive a morphology-dependent response from a
  solved disk PDE before fitting.
\end{enumerate}

% =============================================================================
\section{Outlook: retained results and a reconstruction path}
\label{sec:path}

\subsection{What this note establishes}

\begin{enumerate}
\item The coincidence \eqref{eq:coincidence} is empirical motivation only.
\item \(\azero\equiv c\Hzero/2\pi\) is allowed solely as a
  \emph{present-epoch phenomenological postulate}.
\item The weak-field combination
  \(\Cobs=\Cm^{3/2}/\sqrt{\CIR}\) is a derived, normalization-invariant
  observable coefficient [Eq.~\eqref{eq:Cobs}].
\item On compact \(\Tthree\), sources must be compensated and disk solutions
  generally admit a curl term [Eq.~\eqref{eq:curl}].
\item The no-go results of Secs.~\ref{sec:a0}--\ref{sec:rar} rule out the
  usual cubic geometric packaging of those two numbers.
\end{enumerate}

\subsection{Path required before any geometric-derivation claim}

A subsequent work that claims to \emph{derive} \(\azero\) or \(\Cobs\) from
topology or microphysics must close the following chain.  Items are grouped
by sector; each is a separate technical gate, not a slogan.

\paragraph{Microphysics and circulation.}
\begin{enumerate}
\item[(M1)] Define a condensate order parameter
  \(\Psi=\sqrt{\rho}\,e^{i\theta}\)
  (or a relativistic analogue) with a definite mass parameter \(m\).
\item[(M2)] Obtain the circulation from \(\nabla\theta\) and impose integer
  winding on non-contractible cycles of \(\Tthree\), yielding
  \(\Gamma=n\,h/m\) rather than \(\Gamma=cL\).
\item[(M3)] Fix a single geometric meaning of \(L\)
  (cube edge, closed geodesic length, or circumference) and use it
  consistently in every formula.
\item[(M4)] Derive a map
  \((\Gamma,m,n,L)\mapsto\azero\)
  that does not rely on a chart-dependent ``distribution'' of \(c\Hzero\) over
  an angular range of \(2\pi\).
\end{enumerate}

\paragraph{Cosmic scale and moduli.}
\begin{enumerate}
\item[(S1)] Either restrict scale matching to the present epoch, or supply a
  modulus action for \(L_{\mathrm{phys}}(t)\) that does not force the
  coasting condition \(q=0\).
\item[(S2)] Confront the resulting \(L\) with Planck/COMPACT bounds for cubic
  \(E_1\) and for twisted flat alternatives \cite{planck_topology,compact_promise}.
\end{enumerate}

\paragraph{Weak-field coefficient and data.}
\begin{enumerate}
\item[(E1)] If a projector \(h_{ij}\) is essential, insert it into the action,
  specify the dynamics of \(n^{i}\), and compute
  \(\Cobs=\Cm^{3/2}/\sqrt{\CIR}\)
  as a definite number independent of field normalization.
\item[(E2)] Reconcile \(\Cobs^{2}\azero\) with the empirical RAR/BTFR
  normalization (avoid the factor-\(4/9\) conflict of Sec.~\ref{sec:rar}).
\item[(E3)] Solve the compensated periodic nonlinear equation for at least one
  spherical and one disk baryonic source, including any curl contribution.
\end{enumerate}

\paragraph{Representation (already satisfied here).}
\begin{enumerate}
\item[(R1)] Use a face-identified parallelepiped for \(\Tthree\)
  (Fig.~\ref{fig:t3}), never an embedded doughnut \(\Ttwo\).
\end{enumerate}

Until (M1)--(M4), (S1)--(S2), and (E1)--(E3) are closed, the defensible
scientific product remains exactly what this note provides: present-epoch
phenomenological matching of Eq.~\eqref{eq:coincidence}, the derived invariant
\(\Cobs\), and the no-go results for the cubic geometric story---not a
geometric derivation of the MOND scale.

% =============================================================================
\section*{Code and data availability}

Figure~\ref{fig:t3} is produced by
\texttt{Scripts/itsm\_t3\_fundamental\_domain.py} in the public archive
\url{https://github.com/brendohxd/ITSM-Integrated-Toroidal-Syntropic-Model}
(Zenodo DOI~\href{https://doi.org/10.5281/zenodo.18808348}{10.5281/zenodo.18808348}).
Scripts that render an embedded doughnut surface depict \(\Ttwo\) and should
not be used as \(\Tthree\) schematics.  No new observational catalogue is
analysed.

\begin{acknowledgments}
The author thanks colleagues who pressed for explicit claim status and for
separation of postulates from derivations.  Remaining errors are the author's.
\end{acknowledgments}

\bibliography{references}

\end{document}
